\ficha{A voz própria dos jornais: registro para o estudo da imprensa}

\bloco{Fontes lidas}

Textos de redação de \textit{O Paiz} (24 e 25 de setembro e 7 de novembro de
1906; 21 de março de 1907; 1º de dezembro de 1909; 6 e 18 de agosto de 1914) e
do \textit{Correio da Manhã} (21 de setembro e 10 de outubro de 1906), além do
testemunho de David Campista (5 de outubro de 1906) e de Pinheiro Machado
(dezembro de 1913) sobre a imprensa de 1906.

\bloco{Estatuto desta ficha}

As fichas anteriores registram posições de atores. Esta registra o que os
próprios jornais escreveram quando falaram por si, em artigo sem assinatura,
nota de redação ou coluna da casa. É inventário de voz com citação. Não mede a
posição editorial de nenhum jornal: não li as edições que não nomeiam a Caixa,
não pesei frequência e não sei quem escrevia cada texto. A medida depende de
instrumento que ainda não escolhi.

\bloco{O Paiz em 1906}

Listei as 55 primeiras páginas de \textit{O Paiz} de 1906 que nomeiam a Caixa.
De março a agosto só encontrei voz hospedada: entrevistas, textos explicativos
sobre o Convênio de Taubaté e relatos de sessão. Entre 24 de setembro e 10 de
novembro, catorze edições trazem na primeira página trecho hostil ao projeto.
Li três delas em janela longa e conferi duas na imagem.

O artigo de 24 de setembro se chama \textit{Caixa de surpresas} e não tem
assinatura (título e ausência de assinatura conferidos na imagem). Defende o
plano de 1899 e descreve a Caixa como invenção com que Campista pretende

\cit[I:per178691_1906_08026:1]{augmentar o numero das nossas inutilidades
funestas}{\jor{opaiz19060924}{O Paiz}{24 set. 1906}{1}}

O artigo de 7 de novembro, \textit{Um mecanismo simples}, abre a primeira
página (título e posição conferidos na imagem) e responde ao parecer de Urbano
Santos no Senado com uma leitura distributiva da taxa:

\cit[I:per178691_1906_08070:1]{os lavradores querem cambio baixo para ganhar o
agio do ouro}{\jor{opaiz19061107}{O Paiz}{7 nov. 1906}{1}}

A edição de 25 de setembro ataca a expressão \textit{embaraços theoricos}, com
que Campista se referira à promessa legal de 1846
(\jor{opaiz19060925}{O Paiz}{25 set. 1906}{1}). A mesma expressão volta na nota
de 1º de dezembro de 1909 contra a proposta de Campista sobre o depósito (ficha
4), o que sugere a mesma pena, sem que eu possa prová-lo.

Dois testemunhos externos vão no mesmo sentido. Campista, na tribuna, em 5 de
outubro de 1906, diz ter encontrado em \textit{O Paiz}, naquele mesmo ano, as
melhores defesas de suas ideias, e que o jornal passara a combater a reforma
(ficha 2). Pinheiro Machado, em 1913, inclui \textit{O Paiz} entre os órgãos
que combateram a Caixa (ficha 3), e é o próprio jornal que o publica.

\bloco{O Correio da Manhã em 1906}

O texto de primeira página de 21 de setembro escreve contra o projeto e contra
o Convênio, que chama de \tr[I:per089842_1906_01897:1]{aventura sem
precedentes}, e a favor do governo e do Banco do Brasil
(\jor{cm19060921}{Correio da Manhã}{21 set. 1906}{1}). A nota de 10 de outubro
sobre os ex-ministros da Fazenda conclui que só um dos dez é inequivocamente
favorável ao projeto (\jor{cm19061010}{Correio da Manhã}{10 out. 1906}{1}). Campista
lembra que o mesmo jornal defendera a fixação do câmbio em 13 de janeiro de
1906 (ficha 2).

\bloco{O Paiz depois de 1906}

Em 1907, uma coluna satírica acusa o Banco do Brasil de comprimir o câmbio em
15 (ficha 4). Em 1909, a nota de redação defende que o depósito da Caixa
responde pelos bilhetes emitidos (ficha 4). Em 1914, o jornal apoia a emissão
de papel-moeda contra o voto de Antônio Carlos e censura Rui Barbosa por
combater o decreto do feriado (ficha 7).

\bloco{O que se pode dizer}

\begin{enumerate}[leftmargin=*,nosep]
\item Entre o fim de setembro e novembro de 1906, a primeira página de
\textit{O Paiz} publicou uma série de textos contra a Caixa de Campista, em
nome do plano de 1899, da promessa de 1846 e do custo de vida. Dois são artigos
sem assinatura, conferidos na imagem; os outros onze eu vi só pelo trecho em
torno do nome da Caixa, sem verificar título nem assinatura.
\item Em setembro e outubro de 1906 o \textit{Correio da Manhã} escreveu, com
voz própria, contra o projeto e a favor do governo.
\item Os dois registros corroboram, para esses dois jornais, a lista de Pinheiro
Machado. O terceiro nome da lista, o \textit{Jornal do Commercio}, não integra
o corpus.
\end{enumerate}

\bloco{O que não se pode dizer}

Não se pode afirmar, com este inventário, qual foi a posição de \textit{O Paiz}
em 1906. Falta saber quem escrevia a série, se a linha vale para o ano inteiro
(Campista diz que não, e de março a agosto só achei voz hospedada) e qual o
peso desses textos no conjunto das edições. A mesma cautela vale para o
\textit{Correio da Manhã}, que em janeiro de 1906, segundo Campista, defendia a
fixação.
